\chapter*{Introduction Générale}
%ajouter le titre introduction générale dans la table des matières
\addcontentsline{toc}{chapter}{Introduction Générale}
%définir l'entête de ce chapitre 
\markboth{Introduction Générale}{}
\setlength{\intextsep}{6pt}
\setlength{\textfloatsep}{8pt}
\setlength{\floatsep}{6pt}
\setlength{\abovecaptionskip}{3pt}
\setlength{\belowcaptionskip}{1pt}
\setlength{\LTpre}{3pt}
\setlength{\LTpost}{3pt}

À l’ère de la transformation numérique, la gestion des ressources humaines s’impose comme l’un des domaines les plus stratégiques au sein des organisations modernes. Couvrant des processus essentiels tels que le suivi des présences, la gestion des congés, l’administration des salaires et le traitement des documents administratifs, elle influence directement la performance opérationnelle ainsi que le bien-être des collaborateurs. Toutefois, lorsque ces processus reposent encore sur des méthodes manuelles ou des outils traditionnels, ils engendrent souvent des erreurs, des pertes de temps et un manque de transparence pouvant nuire au bon fonctionnement de l’entreprise. Dans ce contexte, l’intégration des technologies intelligentes, notamment la reconnaissance faciale, représente une solution innovante permettant d’automatiser les tâches, d’améliorer la fiabilité des traitements et d’optimiser la gestion du personnel. La digitalisation des processus RH devient ainsi une nécessité incontournable pour les entreprises souhaitant améliorer leur efficacité et renforcer leur compétitivité.

Dans ce cadre, nous avons eu l’opportunité de réaliser notre Projet de Fin d’Études au sein de la société UNIBRO Distribution, entreprise spécialisée dans la distribution de produits alimentaires et de grande consommation. Notre mission a consisté à concevoir et développer un Logiciel RH Intelligent de Pointage Automatique, visant à moderniser et centraliser la gestion des ressources humaines au sein de l’entreprise. Le fonctionnement de cette application repose sur une plateforme web sécurisée accessible selon trois profils utilisateurs : Administrateur, Responsable RH et Employé. Après authentification, chaque utilisateur accède à un espace dédié contenant les fonctionnalités correspondant à son rôle. Le système intègre un module de reconnaissance faciale permettant d’identifier automatiquement les employés via des caméras connectées à la plateforme. Lorsqu’un employé se présente devant la caméra, son visage est analysé puis comparé aux données biométriques enregistrées afin de valider automatiquement son pointage d’entrée ou de sortie en temps réel. En complément du pointage intelligent, l’application offre un ensemble de services RH tels que la gestion des congés, des absences, des primes, des déductions, des salaires et des documents administratifs. Les responsables RH disposent de tableaux de bord interactifs leur permettant de superviser les données du personnel, de générer des rapports et de suivre les indicateurs de présence, tandis que les employés peuvent consulter leurs informations personnelles et effectuer leurs demandes de manière autonome. Afin d’assurer performance, sécurité et évolutivité, la solution a été développée à l’aide des technologies Next.js, Flask et MySQL, tout en suivant la méthodologie agile Scrum pour garantir une organisation efficace du projet.

Ce rapport est organisé en quatre chapitres principaux. Le premier chapitre donne un aperçu général du projet, de l’organisation d’accueil, de l’étude de l’existant et de la méthodologie du projet. Le deuxième chapitre est consacré à l’analyse des besoins du système, à sa description fonctionnelle et non fonctionnelle, à la définition de son architecture et de sa stratégie de réalisation suivant le processus Scrum. Le troisième chapitre aborde la première release du projet, qui comprend la gestion des accès, la gestion des comptes et la réalisation de la fonctionnalité de pointage par reconnaissance faciale. Enfin, le quatrième chapitre concerne la deuxième release du projet, consacrée aux fonctionnalités avancées de gestion de personnel, aux tableaux de bord et aux reportings du projet.
\clearpage
\mtcaddchapter